% LSPV-2 REPORT::UGENCOMP1::2026-10-12T10:22:04Z
% !TEX program = xelatex
\documentclass[11pt]{article}
\usepackage[T1]{fontenc}
\usepackage[margin=0.84in]{geometry}
\usepackage{amsmath,amssymb,amsthm,booktabs,array,xurl,hyperref}
\setlength{\emergencystretch}{3em}
\newtheorem{theorem}{Theorem}
\newtheorem{corollary}[theorem]{Corollary}
\newcommand{\SAT}{\mathsf{SAT}}
\newcommand{\Pclass}{\mathsf{P}}
\newcommand{\PROVED}{\textbf{PROVED/EXACT}}
\newcommand{\KNOWN}{\textbf{KNOWN THEOREM}}
\newcommand{\COMP}{\textbf{COMPUTATIONAL}}
\newcommand{\OPEN}{\textbf{OPEN}}
\newcommand{\RET}{\textbf{FALSIFIED/RETIRED}}

\title{LSPV-2 UGENCOMP1: Charged Uniform Per-Length Compilation Is Exactly Uniform Decision}
\author{Universal Core Tests}
\date{12 October 2026, 10:22 UTC}

\begin{document}
\maketitle

\begin{abstract}
POINTINF1 showed that minimizing independently over a new finite machine at
each input length collapses by exponential hardwiring.  We test the strongest
obvious repair: require one uniform generator for the length-specific code and
charge generation, code output, access, simulation and the final decision.
The repair is exact but not intermediate.  A polynomial-time generator and a
polynomial-time evaluator compose into one polynomial-time global decider;
conversely every polynomial-time global decider supplies such a family using
constant code.  Thus \texttt{CHARGED-UNIFORM-FAMILY-EQUIV1} is
\PROVED, while \texttt{UNIFORM-PER-LENGTH-COMPILATION-AS-WEAKER-BRIDGE1}
is \RET.  For SAT the repaired condition is equivalent to \(P=NP\), not a
new route to it.  A bounded audit made 1,134,422 assertions with zero
failures.  P versus NP remains \OPEN.
\end{abstract}

\section{Target, encoding and exact quantifiers}

Let \(L\subseteq\{0,1\}^*\) be a decision language in the ordinary uniform
deterministic multitape model with read-only input and sequential explicit
output.  A \emph{charged uniform per-length family} is one fixed pair
\((G,U)\) such that, for every \(n\), the generator \(G(1^n)\) emits a finite
code \(z_n\), and
\[
  \forall n\;\forall x\in\{0,1\}^n:\qquad
  U(z_n,x)=L(x).                                      \tag{1}
\]
The same \(G\) and \(U\) serve all lengths.  Let \(g(n),h(n)\) be generation
time and workspace, \(\ell(n)=|z_n|\), and \(t(n),s(n)\) the time and
workspace of the declared evaluation on \((z_n,x)\).  Generation, every code
write, code storage or virtual access, universal simulation, input scan and
one-bit output are charged.  In particular, sequential explicit output gives
\[
  \ell(n)\le g(n)+O(1).                               \tag{2}
\]

The candidate under test was whether polynomial \(g,h,\ell,t,s\) repair the
quantifier failure of POINTINF1 while leaving a weaker obstruction than the
original uniform decision problem.  It does repair the quantifiers, but the
claimed intermediate status is false.

\section{Exact composition theorem}

\begin{theorem}[CHARGED-UNIFORM-FAMILY-EQUIV1]
For every decidable language \(L\), the following are equivalent.
\begin{enumerate}
\item \(L\in\Pclass\).
\item There is one charged uniform per-length family \((G,U)\) satisfying
(1), with fixed polynomial bounds on generation and evaluation including all
code-access costs.
\end{enumerate}
More precisely, the composed machine \(D\) has
\[
 T_D(n)\le O(n)+g(n)+t(n)+O(1),                       \tag{3}
\]
when \(t\) already includes access/simulation, and under explicit retained
code
\[
 S_D(n)\le O(\log n)+h(n)+\ell(n)+s(n)+O(1).          \tag{4}
\]
If code is streamed, recomputed or accessed virtually, the corresponding
cost must instead occur in \(g,h,t,s\); it is not free.
\end{theorem}

\begin{proof}
Assume (2) of the theorem.  Define one ordinary machine \(D\).  On input
\(x\), it computes \(n=|x|\), runs the fixed generator \(G(1^n)\), and then
runs the fixed evaluator \(U(z_n,x)\).  Correctness is (1).  Equations
(3)--(4) follow by literal sequential composition.  Since all component
bounds have fixed polynomial exponents and (2) bounds the materialized code
length, \(D\) is a uniform polynomial-time decider for \(L\).

Conversely, let one fixed machine \(B\) decide \(L\) in polynomial time.
Choose a fixed interpreter \(U\), let \(z\) be the constant code of \(B\),
and let \(G(1^n)=z\) for every \(n\).  Universal simulation has a fixed
polynomial overhead, so the resulting generator/evaluator family satisfies
(1) with polynomial charged bounds.  No nonuniform advice is used. \qed
\end{proof}

\begin{corollary}[SAT endpoint equivalence]
A charged uniform per-length polynomial family for SAT exists if and only if
\(P=NP\).  Consequently, assuming such a family to prove \(P\ne NP\) assumes
the contradiction hypothesis in a different syntax; disproving it for SAT is
already the desired separation.
\end{corollary}

\paragraph{SAME\_REALIZATION.}
The theorem uses one composed machine \(D\) at every length.  Generator,
evaluator, code and output are parts of that same execution.  This is exactly
the missing cross-length realization in POINTINF1.  Once imposed and fully
charged, however, it is simply uniform computation.

\section{Explicit-code/advice trilemma}

POINTINF1's hardwired depth-\(n\) truth-table machine uses
\(2^{n+1}-1\) prefix states and contains \(2^n\) answer bits in its finite
description.  A purported repair has only three possibilities.
\begin{enumerate}
\item \textbf{Explicit code.}  Emitting \(2^n\) table bits costs at least
\(2^n\) sequential writes.  The hidden hardwiring cost reappears.
\item \textbf{Succinct generated code.}  If both generation and evaluation,
including every access and simulation step, are polynomial, the composition
theorem gives a single polynomial decider.
\item \textbf{Opaque external table/oracle.}  If a length-indexed object is
supplied without generation or access cost, the model has changed to advice
or oracle computation.  It is not an ordinary uniform SAT decider.
\end{enumerate}
This proves \texttt{GENERATOR-CHARGE-TRILEMMA1} as an exact accounting
classification.  It does not establish which case SAT occupies.

\section{Bounded reproducible audit}

The deterministic script
\path{experiments/review/uniform_generator_composition_audit.py} exhausts all
65,814 Boolean languages on lengths \(0\) through \(4\), covering 1,050,698
input evaluations.  For explicit tables it counts one charged generator bit
write and one evaluator table read per input word and checks the exact prefix
tree size.  It separately checks all \(7^4=2,401\) generation/code/evaluation
exponent ledgers on seven finite sizes and all \(7^3=343\) retained-workspace
ledgers on three sizes.  The audit made 1,134,422 assertions with zero
failures in 0.264 seconds; no chunking was needed.

Raw output is
\path{results/raw/review/2026-10-12_uniform_generator_composition.json}.
The script SHA-256 is
\texttt{229af0187f4b1b226c4ff56e72d185b77874be23039bfdb3bd7d1cb998d282a7}.
This evidence is \COMP.  It does not prove the asymptotic theorem, benchmark a
SAT solver or justify extrapolation.  The theorem is proved by composition.

Preprocessing is exactly the charged run of \(G\) and its writes.  Extraction
and witness reconstruction are absent.  Evaluation includes input scanning,
code access and simulation.  The output is one bit.  Retained code counts as
workspace; streamed or virtual code moves the same obligation into access or
recomputation cost.

\section{Hostile controls and semantic cost comparison}

The theorem is language-general, so 2SAT, Horn, affine/Tseitin with Gaussian
elimination, bounded treewidth, the signed clean core, independent products,
padding, redundant clauses and alternate encodings do not affect the
composition.  Easy subclasses may have compact polynomial generators and
evaluators; that is an algorithmic shortcut, not evidence of a universal
compatibility bottleneck.  DIRECT/COPY/DEAD/DELAY/INTERP/RECOMP wrappers can
inflate a chosen implementation but cannot invalidate the existence of the
composed decider.  A one-bit shell is already the declared output contract.

Thus semantic table size, number of satisfying assignments, compatibility
states, code Schur-square dimension or truth-table entropy cannot be converted
to runtime without a theorem forcing the same ordinary computation to acquire
or access that representation.  The present result only says that when an
effective polynomial acquisition procedure is supplied, it is itself part of
a polynomial decider.

\section{Primary-literature and barrier-transfer preflight}

The two blocking obstacles were (i) cross-length uniform selection and
(ii) a model-matched unconditional landing after all generation/access costs
are charged.  Primary sources were checked before the new theorem.

\begin{itemize}
\item Karp--Lipton's advice-machine formalism is the directly relevant
uniform/nonuniform distinction.  It is \KNOWN context; no collapse theorem is
used here.  An ungenerated \(z_n\) is advice, while a polynomially generated
and evaluated \(z_n\) composes into a uniform algorithm.
\item Ruzzo's uniform-circuit framework is a directly relevant structural
analogy: a uniform family needs an effective cross-length description
mechanism.  Its alternating-machine characterizations do not prove the
present ordinary-decider theorem; importing circuit depth or gate lower bounds
would require a new model transfer.
\item Fortnow--Lipton--van Melkebeek--Viglas give \KNOWN model-matched SAT
time--space lower bounds for one general machine.  They do not forbid a
polynomial generator/evaluator family without the composition above, and the
composition does not land in their forbidden region.  That landing remains
\OPEN.
\end{itemize}

The electronics analogy of fabricating one chip per width is useful only when
fabrication description, area and access are charged.  Treating chips as free
ordinary-TM finite control is a misleading analogy / failed transfer.  The
exact transferable fact is merely sequential output length lower-bounds
generator time.  VLSI area--time tradeoffs, communication cuts, proof-system
interpolation and pebbling remain useful analogies requiring separate
all-decider compilers.

\section{Barrier audit}

\paragraph{Relativization.}
The composition proof relativizes: if the fixed generator/evaluator use the
same oracle, their composition does too.  Therefore it is a no-go theorem, not
a nonrelativizing hinge.

\paragraph{Algebrization.}
The proof is representation-agnostic and survives algebraic oracle access
when all such access is charged in the declared model.  It supplies no
nonalgebrizing step.

\paragraph{Natural Proofs.}
No large constructive circuit property is asserted.  The result neither
invokes nor escapes the Natural Proofs barrier.

\paragraph{Arbitrary exponents.}
The exponents of \(G\) and \(U\) may be arbitrary fixed constants; the
composed exponent is their maximum up to fixed simulation overhead.  No
solver-independent small exponent is obtained.  This is precisely why the
result cannot produce an FLMV contradiction.

\section{Post-quantum technique ledger}

\begin{itemize}
\item Schur squares, shortening/puncturing, quadratic-relation ranks, Hilbert
functions, syzygies and Betti numbers: \textbf{FAILED TRANSFER}.  No code
representation is forced, and exploitable structure may make \(G,U\) easier.
\item Collision entropy: \textbf{EXACT OBJECT/THEOREM TRANSFER} only for a
fixed specified map, via
\[
 |\operatorname{im}f|\ge
 |A|^2/|\{(y,y'):f(y)=f(y')\}|.
\]
It does not supply a uniform generator or a runtime lower bound.
\item LWE-style worst-case/average-case extraction:
\textbf{ANALOGY/CANDIDATE needing a new bridge}.  Any black-box extractor and
all solver calls compose into a separately charged algorithm.
\item Collision extraction/special soundness:
\textbf{ANALOGY/CANDIDATE needing a construction-specific relation}.  It does
not bind codes across lengths.
\item Hash/random-oracle methods: \textbf{FAILED TRANSFER} unless the oracle
or cryptographic assumption is removed unconditionally.
\end{itemize}

\section{CORPUS REVIEWED}

The run surveyed the current GitHub tree and read root \path{main.tex}, actual
\path{main4.tex} and \path{main5.tex}, \path{AGENTS.md}, \path{CODEX.md}, the
pnp-research skill and protocol, and every current task specification.  It
read newest verified MAIN SYNTHESIS-0023 and historical SYNTHESIS-0022; newest
verified specialist reports POINTINF1, CLKSLACK1, JOINTGAUGE1 and CALQNT1;
and relevant older TESTS, red-team, lower-bound and fixed-program reports.

Repository search covered ``uniform generator'', ``per-length'', ``advice'',
``hardwiring'', ``length-indexed table'', ``compiler'' and ``P-uniform''.
The run inspected FIXED-CODE-ERASURE1 in the defended \path{main9.tex} spine,
INSTANCE-HARDWIRE1, BINARYDOF\_UNIFORM1, POINTINF1 and their cited raw data.
It inspected the HEC, VSC-HEC, G-PPC-HEC, CRS-PPC, CCSPC and RCCN causal
components, governance/write-repair notes, and the current experiment/result
ledger.  LSPV-1's predecessor report POINTINF1 was nonblank, had the required
header and sentinel, and was readable from GitHub; no recovery request was
needed.

\section{DUPLICATION / CONTRADICTION AUDIT}

FIXED-CODE-ERASURE1 already says that a length-indexed table is either
generated at charged cost or is nonuniform advice.  POINTINF1 already says
that requiring a polynomial generator would repair its quantifier failure.
Those warnings are adopted, not claimed as new.  The new increment is the
two-way theorem that the repaired polynomial generator/evaluator condition is
exactly membership in \(\Pclass\), with the explicit SAME\_REALIZATION and
time/space equations (3)--(4), plus the exhaustive accounting audit.

No contradiction was found.  POINTINF1's \(\forall n\exists A_n\) machines
need not be uniformly generated.  UGENCOMP1 instead assumes one \(G\) before
all lengths, so composition is legal.  HEC/CRS/CCSPC lower bounds remain in
restricted online models and do not provide this family for arbitrary static
SAT deciders.

\section{WHAT CHANGED RELATIVE TO ROOT MAIN}

Root MAIN already excludes pointwise per-length minimization through the
LSPV-1 ownership region.  This run closes the immediate repair: adding a
fully charged polynomial uniform selector does not yield a weaker semantic
resource; it restates the global polynomial decision problem.  The LSPV-2
region is reconciled with this trilemma.  No executive theorem, clean-core
definition or exponent endpoint is strengthened.  The global status remains
\OPEN.

\section{ADOPTED / REJECTED / RETIRED PRIOR COMPONENTS}

\begin{itemize}
\item \textbf{ADOPTED:} POINTINF1's hardwiring theorem and quantifier guard;
FIXED-CODE-ERASURE1; ordinary uniform encoding; SAME\_REALIZATION; full
generation/access accounting; and FLMV only in its declared model.
\item \textbf{REJECTED as free transfers:} per-width fabricated hardware,
P-uniform circuit depth/size conclusions, online HEC state bounds, proof-system
lower bounds and advice tables as statements about arbitrary SAT machines.
\item \textbf{RETIRED:}
\texttt{UNIFORM-PER-LENGTH-COMPILATION-AS-WEAKER-BRIDGE1}.  A polynomial
charged version is \(L\in\Pclass\); an uncharged version is advice/model
change; an explicit truth-table version pays exponential output.
\end{itemize}

\section{STALE OR OVERSTATED CLAIMS FOUND}

The phrase ``requiring a uniform polynomial-time compiler is stronger still''
is correct but incomplete if presented as a promising intermediate for the
decision language.  With a polynomial evaluator it is exactly the original
uniform polynomial-time decision condition.  Root MAIN did not claim
otherwise, so no root theorem is contradicted.  Any future statement must
distinguish a generator for a restricted representation from a generator whose
output is already polynomially evaluable on every input.

\section{Surviving missing implication}

The highest-value missing arrow is not another selector for per-length SAT
code.  It is a genuinely resource-improving, model-preserving transformation:
from one assumed polynomial SAT decider of arbitrary exponent, construct the
same fully charged ordinary realization in an unconditional forbidden
time--space region, with exponents independent of that source exponent.
Equivalently, a structural generator is useful only if a separately proved
lower bound constrains its generated representation and an all-decider
compiler lands every assumed solver there without already solving SAT as its
premise.  This remains \OPEN.

One load-bearing arrow is OPEN; therefore P versus NP remains OPEN.

\begin{thebibliography}{9}
\bibitem{KarpLipton1982}
R. M. Karp and R. J. Lipton, ``Turing machines that take advice,''
\emph{L'Enseignement Math\'ematique} 28, 191--209, 1982.

\bibitem{Ruzzo1981}
W. L. Ruzzo, ``On uniform circuit complexity,'' \emph{Journal of Computer
and System Sciences} 22(3), 365--383, 1981.
doi:10.1016/0022-0000(81)90038-6.

\bibitem{FLMV2005}
L. Fortnow, R. Lipton, D. van Melkebeek, and A. Viglas,
``Time-space lower bounds for satisfiability,'' \emph{Journal of the ACM}
52(6), 835--865, 2005. doi:10.1145/1101821.1101822.
\end{thebibliography}

% END_OF_VERIFIED_REPORT::LSPV-2::UGENCOMP1
\end{document}
